\documentclass[11pt,a4paper]{article}
\usepackage{curso}
\begin{document}
\portada{Ejercicios --- Lección}{4.8}{RBMs y divergencia contrastiva}
\section{Ejercicios}
\begin{enumerate}[label=\textbf{E\arabic*.}]
\item (B) Derive $p(h_j=1\mid v)=\sigma(c_j+v^\top W_{:j})$.
\item (B) Derive la energía libre $F(v)=-b^\top v-\sum_j\ln(1+e^{c_j+v^\top W_{:j}})$.
\item (B) ¿Por qué CD-$k$ es sesgada y qué ocurre cuando $k\to\infty$?
\item (B) ¿Cuántos estados hay que sumar para $Z$ con 9 visibles (sumando las ocultas analíticamente)?
\item (I) Muestre que $\partial F(v)/\partial W_{ij}=-v_i\,\sigma(c_j+v^\top W_{:j})$ y escriba el gradiente exacto de la log-verosimilitud con $F$.
\item (I) ¿Cuántas fases de reloj usa un paso de CD-10 en hardware de pbits con una capa por fase?
\item (A) Reproduzca las curvas CD-1 y CD-10 y añada PCD con $\eta=0{,}01$ y 20\,000 pasos. ¿Supera a CD-10?
\end{enumerate}
\section{Soluciones}
\begin{solucion}
\textbf{E1.} Los términos con $h_j$ son $h_j(c_j+v^\top W_{:j})$; $p(h_j=1)/p(h_j=0)=e^{c_j+v^\top W_{:j}}$.

\textbf{E2.} $\sum_he^{-E}=e^{b^\top v}\prod_j(1+e^{c_j+v^\top W_{:j}})$.

\textbf{E3.} La cadena no alcanza el equilibrio en $k$ pasos; con $k\to\infty$ el término converge a $\langle vh^\top\rangle_{\text{modelo}}$.

\textbf{E4.} $2^9=512$.

\textbf{E5.} $\nabla\mathcal L=-\overline{\nabla F}_{\text{datos}}+\langle\nabla F\rangle_{p(v)}$.

\textbf{E6.} $2k+2=22$ medias fases (más la lectura de estadísticas).

\textbf{E7.} Con tasa pequeña PCD suele igualar o superar a CD-10 en verosimilitud, a cambio de más pasos.
\end{solucion}
\end{document}
